\documentclass{article}
\usepackage{macros-article}


\begin{document}
\section{Flowcharts}

\setlength{\columnsep}{1cm}
    \begin{multicols}{2}
    \begin{gitFlowCard}{Create local Repo then push to remote}
        \begin{tikzpicture}[
            node distance=0.8cm,
            dot/.style={circle, draw=gitRed, fill=white, line width=1.5pt, inner sep=2pt},
            line/.style={gitRed, line width=1.5pt},
            dashed line/.style={gitRed, line width=1.5pt, dashed}
        ]

            % Main Trunk Nodes
            \node[dot] (n1) {};
            \node[right=2mm of n1, font=\sffamily\bfseries] {Create repo with hosting service};
            
            \node[dot, below=of n1] (n2) {};
            \node[right=2mm of n2] {\gitlabel{git init} \textsf{\small in your local dir}};
            
            \node[dot, below=of n2] (n3) {};
            \node[right=2mm of n3] {\gitlabel{git remote add origin <url>}};

            % The Branching Point
            \node[below=1.5cm of n3] (n4) {}; % Invisible anchor for dashed line
            
            % Side Branch for conditional
            \coordinate (branchStart) at ($(n3)!0.48!(n4)$);
            \node[right=1.2cm of branchStart, text=gitRed, font=\sffamily\scriptsize\itshape] (condText) {If the created repo has commits already (e.g. a README)};
            
            \node[dot, below=0.4cm of condText.south west, anchor=center] (b1) {};
            \node[right=2mm of b1] {\gitlabel{git fetch}};
            
            \node[dot, below=of b1] (b2) {};
            \node[right=2mm of b2] {\gitlabel{git branch -u origin/main main}};
            
            \node[dot, below=of b2] (b3) {};
            \node[right=2mm of b3] {\gitlabel{git pull --rebase}};



            % Resume Main Trunk
            \node[dot, below=4.5cm of n3] (n5) {};
            \node[right=2mm of n5] {\gitlabel{git add}};
        
            \node[dot, below=of n5] (n6) {};
            \node[right=2mm of n6] {\gitlabel{git commit -m "message"}};
            
            \node[dot, below=of n6] (n7) {};
            \node[right=2mm of n7] {\gitlabel{git push -u origin main}};

            \coordinate (branchEnd) at ($(n4)!0.8!(n5)$);
            
            % Drawing the lines
            \draw[line] (n1) -- (n2);
            \draw[line] (n2) -- (n3);
            \draw[dashed line] (n3) -- (n5);
            \draw[line] (n5) -- (n6);
            \draw[line] (n6) -- (n7);
            
            % Drawing the branch curves
            \draw[line, rounded corners=5pt] (branchStart) -- ++(0.5,0) |- (b1);
            \draw[line] (b1) -- (b2);
            \draw[line] (b2) -- (b3);
            \draw[line, rounded corners=5pt] (b3) -- ++(0,-0.35) -| (branchEnd);

        \end{tikzpicture}
    \end{gitFlowCard}

    \begin{gitFlowCard}{Remove and untrack a mistaken file}
        \begin{tikzpicture}[
            node distance=0.8cm,
            dot/.style={circle, draw=gitRed, fill=white, line width=1.5pt, inner sep=2pt},
            line/.style={gitRed, line width=1.5pt}
        ]
            \node[dot] (n1) {};
            \node[right=2mm of n1] {\gitlabel{git rm --cached <file>}};

            \node[dot, below=of n1] (n2) {};
            \node[right=2mm of n2, font=\sffamily\bfseries\small] {add file to your .gitignore};

            \node[dot, below=of n2] (n3) {};
            \node[right=2mm of n3] {\gitlabel{git add .gitignore}};

            \node[dot, below=of n3] (n4) {};
            \node[right=2mm of n4] {\gitlabel{git commit -m "message"}};

            \node[dot, below=of n4] (n5) {};
            \node[right=2mm of n5] {\gitlabel{git push}};

            \draw[line] (n1) -- (n2) -- (n3) -- (n4) -- (n5);

            \node[below=0.2cm of n5.south west, anchor=north west, text=gitRed, font=\sffamily\scriptsize\itshape, text width=5cm] 
                {NOTE: when others pull from remote, <file> will be removed from their local};
        \end{tikzpicture}
    \end{gitFlowCard}
 
    \begin{gitFlowCard}{Discard everything up to last commit}
        \begin{tikzpicture}[
            node distance=0.8cm,
            dot/.style={circle, draw=gitRed, fill=white, line width=1.5pt, inner sep=2pt},
            line/.style={gitRed, line width=1.5pt}
        ]
            \node[dot] (n1) {};
            \node[right=2mm of n1] {\gitlabel{git reset --hard HEAD}};
            
            \node[below=0.4cm of n1.south east, anchor=north west, font=\sffamily\itshape\small] {untracked files remain};

            \node[dot, below=1.2cm of n1] (n2) {};
            \node[right=2mm of n2] {\gitlabel{git clean -fd}};
            
            \node[below=0.4cm of n2.south west, anchor=north west, font=\sffamily\itshape\small] {-d: removes untracked, but not ignored files};

            \draw[line] (n1) -- (n2);
        \end{tikzpicture}
    \end{gitFlowCard}

    \columnbreak

    \begin{gitFlowCard}{Restore file to match last commit}
        \begin{tikzpicture}[
            node distance=0.8cm,
            dot/.style={circle, draw=gitRed, fill=white, line width=1.5pt, inner sep=2pt},
            line/.style={gitRed, line width=1.5pt}
        ]
            \node[dot] (n1) {};
            \node[right=2mm of n1] {\gitlabel{git restore <file>}};
            \node[below=0.4cm of n1.south east, anchor=north west, font=\sffamily\itshape\small] {Reverts \texttt{<file>} in working directory};

            \node[dot, below=1.2cm of n1] (n2) {};
            \node[right=2mm of n2] {\gitlabel{git restore --staged <file>}};
            
            \node[below=0.4cm of n2.south west, anchor=north west, font=\sffamily\itshape\small, align=right] 
                {Removes file from the staging area\\ Keep changes in working directory};

            \draw[line] (n1) -- (n2);
        \end{tikzpicture}
    \end{gitFlowCard}

    \begin{gitFlowCard}{Release new version: GitHub}
        \begin{tikzpicture}[
            node distance=0.8cm,
            dot/.style={circle, draw=gitRed, fill=white, line width=1.5pt, inner sep=2pt},
            grayDot/.style={circle, draw=gray, fill=white, line width=1.5pt, inner sep=2pt},
            line/.style={gitRed, line width=1.5pt},
            grayLine/.style={gray, line width=1.5pt}
        ]
            % Main Trunk
            \node[dot] (n1) {};
            \node[right=2mm of n1, font=\sffamily\bfseries\small] {Make any changes and commit them as normal};

            \node[dot, below=of n1] (n2) {};
            \node[right=2mm of n2] {\gitlabel{git tag -a <version No.> -m "message"}};
            \node[below=0.3cm of n2.south east, anchor=north west, font=\sffamily\itshape\scriptsize] {*-a: annotated tags};

            \node[dot, below=1.1cm of n2] (n3) {};
            \node[right=2mm of n3] {\gitlabel{git push --follow-tags}};

            % Branch off for GitHub UI
            \coordinate (branchPoint) at ($(n3)!0.5!(n3 |- 0,-4.5)$); % Point between dots for curve
            
            \node[grayDot, below right=0.6cm and 1.2cm of n3] (g1) {};
            \node[right=2mm of g1, font=\sffamily\bfseries\small] {Go to Releases $\rightarrow$ Draft new release};

            \node[grayDot, below=of g1] (g2) {};
            \node[right=2mm of g2, font=\sffamily\small] {Choose your tag};

            \node[grayDot, below=of g2] (g3) {};
            \node[right=2mm of g3, font=\sffamily\small] {Target: main};

            \node[grayDot, below=of g3] (g4) {};
            \node[right=2mm of g4, font=\sffamily\small] {Add title and notes};

            \node[grayDot, below=of g4] (g5) {};
            \node[right=2mm of g5, font=\sffamily\small] {Publish};

            % Connections
            \draw[line] (n1) -- (n2);
            \draw[line] (n2) -- (n3);
            \draw[line, rounded corners=8pt] (n3) |- (g1);
            \draw[grayLine] (g1) -- (g2) -- (g3) -- (g4) -- (g5);

        \end{tikzpicture}
    \end{gitFlowCard}

    \begin{gitFlowCard}{Create and Merge branch}
        \begin{tikzpicture}[
            node distance=0.8cm,
            dot/.style={circle, draw=gitRed, fill=white, line width=1.5pt, inner sep=2pt},
            line/.style={gitRed, line width=1.5pt},
            dashedLine/.style={gitRed, line width=1.5pt, dashed}
        ]
            \node[dot] (n1) {};
            \node[right=2mm of n1] {\gitlabel{git switch -c <branch-name>}};

            \node[below=0.4cm of n1, anchor=north west, text=gitRed!60!white, font=\sffamily\scriptsize\itshape] (note1) {do your stuff};

            \node[dot, below=1.2cm of n1] (n2) {};
            \node[right=2mm of n2] {\gitlabel{git switch main}};

            \node[dot, below=of n2] (n3) {};
            \node[right=2mm of n3] {\gitlabel{git merge <branch-name>}};

            \node[below=0.4cm of n3, anchor=north west, text=gitRed!60!white, font=\sffamily\scriptsize\itshape] (note2) {clean up};

            \node[dot, below=1.2cm of n3] (n4) {};
            \node[right=2mm of n4] {\gitlabel{git push origin --delete <branch-name>}};

            \node[dot, below=of n4] (n5) {};
            \node[right=2mm of n5] {\gitlabel{git branch -d <branch-name>}};

            % Drawing lines with dashed segments
            \draw[line] (n1) -- ($(n1)!0.3!(n2)$);
            \draw[dashedLine] ($(n1)!0.3!(n2)$) -- ($(n1)!0.7!(n2)$);
            \draw[line] ($(n1)!0.7!(n2)$) -- (n2);

            \draw[line] (n2) -- (n3);

            \draw[line] (n3) -- ($(n3)!0.3!(n4)$);
            \draw[dashedLine] ($(n3)!0.3!(n4)$) -- ($(n3)!0.7!(n4)$);
            \draw[line] ($(n3)!0.7!(n4)$) -- (n4);

            \draw[line] (n4) -- (n5);
        \end{tikzpicture}
    \end{gitFlowCard}

    \end{multicols}

\newpage

\section{User Settings \& Repo creation}
    \setlength{\columnsep}{1cm}
    \begin{multicols}{2}

        \begin{gitBaseCard}{User Settings}
            \begin{gitsection}{Display gloabal user email \& name}
                \gitCommandLong{???}{Display gloab settings}
                \gitflag{0.62}{git settings user.name}{0.25}{dispplay current repos set git user name}
                \gitflag{0.62}{git settings user.email}{0.25}{dispplay current repos set git user email}
                \gitflag{0.32}{--gloabl}{0.6}{displays or modifies gloab account settings}
            \end{gitsection}
            
            \vspace{2mm}
            
            \begin{gitsection}{Edit .ssh file for more granular user choices}
                ...where to find the .ssh file
            \end{gitsection}
        \end{gitBaseCard}

    \columnbreak

        \begin{gitBaseCard}{Creating Repository}
            \begin{gitsection}{...}
                \gitCommandLong{git clone <url>}{Download an existing repository}
            \end{gitsection}
            
            \vspace{2mm}
            
            \begin{gitsection}{...}
                \gitCommandLong{git init}{Create new local repository}
                \gitflag{0.45}{<project name>}{0.5}{if added: creates .git in sub-dir}
            \end{gitsection}
            
            \vspace{2mm}
            
            \begin{gitsection}{...}
                \gitCommandLong{git remote add origin <url>}{Links the remote repo to your local one}
            \end{gitsection}
        \end{gitBaseCard}

        \begin{gitBaseCard}{Interacting with staging area}
            \begin{gitsection}{...}
                \gitCommandLong{git add <file>}{Add a file in it's current shape to your next commit (stage)}
                \gitflag{0.35}{add .}{0.6}{adds every file to staging area}
                \gitflag{0.35}{add *}{0.6}{CAUTION: don't use this - it is bash syntax}
            \end{gitsection}
            
            \vspace{2mm}
            
            \begin{gitsection}{...}
                \gitCommandLong{git reset <file>}{Unstage a file while retaining the changes in the working directory}
                \gitflag{0.3}{--hard}{0.6}{???}
            \end{gitsection}
        \end{gitBaseCard}
    \end{multicols}

\newpage

\section{Some other title}
    \setlength{\columnsep}{1cm}
    \begin{multicols}{2}

        \begin{gitBaseCard}{Commits}
            \begin{gitsection}{...}
                \gitCommandLong{git commit -m "message"}{Commit all staged files with a commit message}
                \gitflag{0.45}{-am "message"}{0.5}{add all tracked files and then commit}
                \gitflag{0.3}{--amend}{0.6}{amend a commit}
            \end{gitsection}
            
            \vspace{2mm}
            
            \begin{gitsection}{...}
                \gitCommandLong{git reset --soft HEAD~1 }{Reset commit}
                \gitflag{0.45}{--hard <commit>}{0.5}{staged changes vs last commits}
            \end{gitsection}

            \vspace{2mm}

            \begin{gitsection}{...}
                \gitCommandLong{git tag <tag-name> }{Tag your last commit}
                \par\textsf{\small *tags needs to be committed separately}
            \end{gitsection}
        \end{gitBaseCard}

        \begin{gitBaseCard}{Stash}
            \begin{gitsection}{...}
                \gitCommandLong{git stash}{Stash all unstaged files}
                \gitflag{0.3}{list}{0.6}{list all stashed files}
                \gitflag{0.3}{pop}{0.6}{moves files from stash to active working dir}
                \gitflag{0.3}{drop}{0.6}{removes all files from stash}
            \end{gitsection}
        \end{gitBaseCard}

        \columnbreak

        \begin{gitBaseCard}{Branches}
            \begin{gitsection}{...}
                \gitCommandLong{git branch}{List all local branches}
                \gitflag{0.2}{-a}{0.75}{list all local and remote branches}
                \gitflag{0.2}{-av}{0.75}{with extra info}
                \gitflag{0.45}{-m <new-name>}{0.45}{rename (move) current branch to new}
                \gitflag{0.45}{-c <new-name>}{0.45}{copy current branch to new-name}
                \gitflag{0.2}{-d}{0.75}{delete local branch}
                
                \gitCommandLong{git push origin --delete <branch>}{deletes a remote branch}

            \end{gitsection}
            
            \vspace{2mm}
            
            \begin{gitsection}{...}
                \gitCommandLong{git branch -c <name>}{Copy branch into a new branch}
                \gitCommandLong{git switch -c <name>}{Create a new branch and switch to it}
            \end{gitsection}

            \vspace{2mm}
            
            \begin{gitsection}{...}
                \gitCommandLong{git switch <existing name>}{switch branches}
            \end{gitsection}

             \vspace{2mm}
            
            \begin{gitsection}{...}
                \gitCommandLong{git switch --track <origin/branch-name>}{Creates a local branch that tracks a remote branch and switches to it}
            \end{gitsection}
        \end{gitBaseCard}

    \end{multicols}

\newpage

\section{Some other title}
    \setlength{\columnsep}{1cm}
    \begin{multicols}{2}

        \begin{gitBaseCard}{Remote branch interaction}
            \begin{gitsection}{...}
                \gitCommandLong{git push}{Move local commited changes to remote}
                \gitflag{0.45}{--tags}{0.5}{push any commited tags to remote}
            \end{gitsection}
            
            \vspace{2mm}
            
            \begin{gitsection}{...}
                \gitCommandLong{git push -u origin <branch-name>}{Creates a new remote branch and sets your local upstream}
                \gitCommandLong{git push --force-with-lease}{Safe way to force push: will not push if others habe commited changes}
            \end{gitsection}

            \vspace{2mm}

            \begin{gitsection}{...}
                \gitCommandLong{git pull}{Pull from remote to yours local working directory}
                \gitflag{0.45}{--rebase}{0.4}{pull but rebase instead of merge}
                \gitCommandLong{git fetch}{Pull the changes into your local but not working directory}
            \end{gitsection}
        \end{gitBaseCard}

    \columnbreak

        \begin{gitBaseCard}{File manipulation}
            \begin{gitsection}{...}
                \gitCommandLong{git mv <old> <new>}{Move file to different location}
                \gitCommandLong{git rm <file>}{Remove local file}
            \end{gitsection}

            \vspace{2mm}

            \begin{gitsection}{...}
                \gitflag{0.3}{git clean}{0.3}{???}
                \gitflag{0.55}{git restore <file>}{0.4}{restore a deleted file}
                \gitCommandLong{git restore --staged --worktree <file>}{Restore a file to a previous commit version}
            \end{gitsection}
        \end{gitBaseCard}
    \end{multicols}

\newpage

\section{Some other title}
    \setlength{\columnsep}{1cm}
    \begin{multicols}{2}
        \begin{gitBaseCard}{History archeology}
            \begin{gitsection}{...}
                \gitflag{0.4}{git status}{0.5}{}
                \gitCommandLong{git diff}{Shows the content differences between commits, staging area, and working directory.}
                \gitflag{0.3}{--staged}{0.6}{staged changes vs last commits}
                \gitflag{0.3}{HEAD}{0.6}{working directory vs last commit}
            \end{gitsection}

            \vspace{2mm}

            \begin{gitsection}{...}
                \gitflag{0.3}{git log}{0.3}{shows basic log of recent commits}
                \gitCommandLong{git log --pretty=format:"\%h\%x09\%an\%x09\%ad\%x09\%s"}{outputs: commit hash | author | date | commit message}
            \end{gitsection}

            \begin{gitsection}{...}
                \gitflag{0.3}{git show}{0.3}{output: ???}
                \gitflag{0.3}{git blame}{0.3}{output: ???}
                \gitflag{0.35}{git reflog}{0.5}{output: ???}
            \end{gitsection}
        \end{gitBaseCard}

        \columnbreak

        \begin{gitBaseCard}{Combine histories}
            \begin{gitsection}{...}
                \gitCommandLong{git merge <branch>}{merge <branch> into current branch}
                \gitCommandLong{git rebase <branch>}{Apply any commits of current branch ontop of <branch>???}
            \end{gitsection}
        \end{gitBaseCard}

        \begin{gitBaseCard}{Tags}
            \begin{gitsection}{...}
                \gitCommandLong{git tag <tag-name>}{add a tag to current commit}
                \gitCommandLong{git push --tags}{tags must be pushed separately from other commits}
            \end{gitsection}
        \end{gitBaseCard}

    \end{multicols}

\end{document}